\section*{Bab \ref{ch:g12:comb} --- Kombinatorika dan Pencacahan}

\begin{solution}{exo:g12:comb:1}
Asas perkalian:
$26^2 \times 10^3 \times 26^2 = 26^4 \times 10^3 = 456\,976\,000$.

Tanpa karakter yang berulang, keempat hurufnya harus berbeda
($26 \times 25 \times 24 \times 23$ cara, dengan mengisi kedudukan hurufnya
berurutan) dan ketiga angkanya berbeda ($10 \times 9 \times 8$):
\[
26 \times 25 \times 24 \times 23 \times 10 \times 9 \times 8
= 358\,800 \times 720 = 258\,336\,000 .
\]
\end{solution}

\begin{solution}{exo:g12:comb:2}
$\dbinom83 = \dfrac{8 \times 7 \times 6}{3!} = 56$;
$\dbinom{10}{8} = \dbinom{10}{2} = \dfrac{10 \times 9}{2} = 45$;
\[
\frac{\binom n2}{\binom{n+1}2}
= \frac{n(n-1)/2}{(n+1)n/2} = \frac{n-1}{n+1}.
\]
\end{solution}

\begin{solution}{exo:g12:comb:3}
Pilihlah panitianya: $\binom{30}{4}$ cara. Lalu pilihlah ketua dan
bendaharanya di antara yang 4 itu, berurutan: $4 \times 3 = 12$ cara. Seluruhnya
\[
\binom{30}{4} \times 12 = 27\,405 \times 12 = 328\,860 .
\]
\end{solution}

\begin{solution}{exo:g12:comb:4}
\[
(x+2)^5 = x^5 + 10x^4 + 40x^3 + 80x^2 + 80x + 32 ,
\]
\[
(1-x)^6 = 1 - 6x + 15x^2 - 20x^3 + 15x^4 - 6x^5 + x^6 .
\]
Pada $(2x+3)^7$, suku dalam $x^3$ adalah
$\binom{7}{3}(2x)^3\,3^4 = 35 \times 8 \times 81\, x^3$, sehingga koefisiennya
$22\,680$.
\end{solution}

\begin{solution}{exo:g12:comb:5}
\emph{1.} $\dbinom{52}{5} = 2\,598\,960$.

\emph{2.} Tepat satu kartu raja: pilihlah kartunya ($4$ cara) lalu lengkapi dengan $4$
kartu bukan raja: $4 \times \binom{48}{4} = 4 \times 194\,580 = 778\,320$.
Sekurang-kurangnya satu kartu raja, lewat pencacahan \omterm{def:g10:proba:operations}{komplemennya}:
$\binom{52}{5} - \binom{48}{5} = 2\,598\,960 - 1\,712\,304 = 886\,656$.

\emph{3.} Pilihlah peringkat bagi tiga kartu sekawannya ($13$), warnanya
($\binom43 = 4$), peringkat bagi pasangannya ($12$ yang tersisa), lalu warnanya
($\binom42 = 6$):
$13 \times 4 \times 12 \times 6 = 3744$.
\end{solution}

\begin{solution}{exo:g12:comb:6}
Kata MATH mempunyai 4 huruf yang berbeda: $4! = 24$ anagram.

Kata BANANA mempunyai 6 huruf: tiga A, dua N, satu B. Pilihlah kedudukan
huruf A-nya ($\binom63$), lalu kedudukan huruf N di antara sisanya ($\binom32$),
dan huruf B mengambil tempat terakhirnya:
\[
\binom{6}{3}\binom{3}{2} = 20 \times 3 = 60 .
\]
(Setara dengan $\frac{6!}{3!\,2!\,1!} = 60$.)
\end{solution}

\begin{solution}{exo:g12:comb:7}
\emph{Secara aljabar:}
\[
k\binom nk = \frac{k\,n!}{k!(n-k)!} = \frac{n!}{(k-1)!\,(n-k)!}
= n\,\frac{(n-1)!}{(k-1)!\bigl((n-1)-(k-1)\bigr)!} = n\binom{n-1}{k-1}.
\]
\emph{Lewat pencacahan ganda:} cacahlah pasangan (panitia beranggota $k$, ketuanya).
Pilihlah panitianya dahulu ($\binom nk$) lalu ketuanya ($k$): ada
$k\binom nk$ pasangan. Atau pilihlah ketuanya dahulu ($n$ pilihan) lalu
$k-1$ anggota lainnya dari $n-1$ yang tersisa: ada $n\binom{n-1}{k-1}$ pasangan.
\end{solution}

\begin{solution}{exo:g12:comb:8}
Sebuah lintasan terdiri atas tepat $m + n$ langkah, dengan $m$ langkah ke timur dan
$n$ ke utara; lintasan itu sepenuhnya ditentukan oleh himpunan saat (di antara
yang $m+n$) ketika langkahnya ke timur. Ada $\binom{m+n}{m}$ pilihan semacam itu.
\end{solution}

\begin{solution}{exo:g12:comb:9}
Belahlah himpunan berisi $m + n$ orang menjadi kelompok $A$ beranggota $m$ dan
kelompok $B$ beranggota $n$. Suatu himpunan bagian beranggota $k$ memuat sejumlah
$j$ anggota $A$ ($0 \leq j \leq k$) dan $k - j$ anggota $B$; untuk $j$ yang tetap
ada $\binom mj \binom{n}{k-j}$ himpunan bagian semacam itu, dan asas penjumlahan
atas $j$ memberi identitas Vandermonde.

Dengan $m = n = k$:
\[
\binom{2n}{n} = \sum_{j=0}^n \binom nj \binom{n}{n-j}
= \sum_{j=0}^n \binom nj^{2},
\]
dengan memakai kesetangkupan $\binom{n}{n-j} = \binom nj$.
\end{solution}

\begin{solution}{exo:g12:comb:10}
\emph{Lewat \cref{exo:g12:comb:7}:}
\[
\sum_{k=1}^n k\binom nk = \sum_{k=1}^n n\binom{n-1}{k-1}
= n\sum_{j=0}^{n-1}\binom{n-1}{j} = n\,2^{n-1},
\]
menurut \cref{cor:g12:comb:sums}. \emph{Lewat penurunan:} menurunkan
$(1+x)^n = \sum_k \binom nk x^k$ memberi
$n(1+x)^{n-1} = \sum_k k \binom nk x^{k-1}$; lalu hitung nilainya di $x = 1$.
\end{solution}

\begin{solution}{pb:g12:comb:1}
\textbf{1.} Ada $26^2 \times 10^3 = 676\,000$ pelat. Kata BANANA: $6$
huruf dengan A yang tiga kali dan N yang dua kali, jadi
$\frac{6!}{3!\,2!} = 60$ anagram.

\textbf{2.} Ada $\binom{32}{5} = 201\,376$ tangan; dan
$\binom42 \binom{28}{3} = 6 \times 3\,276 = 19\,656$ tangan dengan
tepat dua kartu raja.

\textbf{3.} Sebuah lintasan adalah kata dengan $4$ huruf R dan $3$ huruf U:
pilihlah kedudukan huruf U-nya, jadi $\binom73 = 35$.

\textbf{4.} $(1+x)^4 = 1 + 4x + 6x^2 + 4x^3 + x^4$. Di
$x = 1$: $\sum_k \binom nk = 2^n$; di $x = -1$:
$\sum_k (-1)^k \binom nk = 0$ --- yaitu jumlah barisnya dan jumlah baris
berselang-seling pada \omterm{prop:g11:binom:pascal}{segitiga Pascal}.

\textbf{5.} Panitia beranggota $k$ orang beserta ketuanya, dari $n$ orang:
pilihlah panitianya lalu ketuanya ($\binom nk \times k$), atau
ketuanya lalu anggota lainnya
($n \times \binom{n-1}{k-1}$): keduanya sama. Dengan menjumlahkan atas $k$,
ruas kanannya berjumlah $n \sum_j \binom{n-1}{j} = n\,2^{n-1}$.

\textbf{6.} Deretan $10$ bintang dan $3$ batang menyandikan pesanannya
(sendok rasa 1 sebelum batang pertama, dan seterusnya); deretannya berisi
$13$ lambang dan ditentukan oleh kedudukan batangnya, jadi ada
$\binom{13}{3} = 286$ pesanan.

\textbf{7.} Ada $12$ bintang dan $2$ batang, jadi $\binom{14}{2} = 91$.

\textbf{8.} Dengan $x', y', z' \geq 0$ dan
$x' + y' + z' = 9$: ada $\binom{11}{2} = 55$.

\textbf{9.} Sebuah monomial $a^i b^j c^k$ dengan $i + j + k = 5$:
ada $\binom72 = 21$.

\textbf{10.} Menurut rumusnya: $3$ bintang, $1$ batang, jadi $\binom41 = 4$;
dan daftarnya: $(3,0)$, $(2,1)$, $(1,2)$, $(0,3)$, jadi cocok.

\textbf{11.} ``Identik'' masuk ketika sebuah pesanan dinyatakan
tak lebih daripada \emph{cacah} tiap rasanya --- bintangnya
tidak membawa nama. Jika sendoknya dimakan berurutan, masing-masing dari
$10$ kedudukan yang berbeda itu memilih rasanya secara bebas, sehingga ada
$4^{10} = 1\,048\,576$ \omterm{def:g12:seq:sequence}{barisan} --- model yang lain dan dunia
yang lain (\cref{met:g12:comb:model}: selalu tanyakan
\emph{berurutan? berbeda? boleh berulang?}).

\textbf{12.} $D_1 = 0$; $D_2 = 1$ (bertukar); $D_3 = 2$ (dua
kitaran-$3$); $D_4 = 9$.

\textbf{13.} Tamu 1 menerima topi $k \neq 1$: ada $n - 1$ pilihan. Jika
tamu $k$ menerima topi 1, maka $n - 2$ tamu sisanya menukar seluruh
topinya sendiri: ada $D_{n-2}$ cara. Jika tamu $k$ \emph{tidak}
menerima topi 1, namailah ulang topi 1 sebagai topi terlarang bagi tamu $k$:
maka $n - 1$ tamu sisanya tertukar semua, ada $D_{n-1}$ cara. Jadi
$D_n = (n-1)(D_{n-1} + D_{n-2})$. Periksalah:
$D_4 = 3(2 + 1) = 9$; dan $D_5 = 4(9 + 2) = 44$.

\textbf{14.} Dari $3! = 6$ pemberian itu, kurangkan yang
menetapkan sekurang-kurangnya satu topi: ada tiga yang menetapkan satu topi
tertentu ($2!$ masing-masing, $3 \times 2 = 6$), yang melebihkan cacah
pasangannya ($3$ pasangan, $1!$ masing-masing) sehingga harus dikembalikan,
lalu identitasnya dikurangkan lagi ($1$): jadi $D_3 = 6 - 6 + 3 - 1 = 2$, yaitu
$3!\left(1 - 1 + \frac12 - \frac16\right) = 2$. Secara umum
$D_n = n!\sum_{k=0}^{n} \frac{(-1)^k}{k!}$.

\textbf{15.} $\frac{D_5}{120} = \frac{44}{120} \approx
0.3667$, sudah dekat dengan $\frac1\eu \approx 0.3679$: jumlah
berselang-seling $1 - 1 + \frac{1}{2!} - \frac{1}{3!} + \dots$
berbaris menuju $\eu^{-1}$. Topi pada pesta besar tertukar semua kira-kira
$36.8\,\%$ waktunya --- tetapan milik undian dan sekretaris
itu, untuk penampakan ketiganya.

\textbf{16.} $\P(\text{sah}) = \frac{D_{10}}{10!} \approx
0.368$. Setiap pengambilan ulang berhasil dengan peluang
$\approx \frac1\eu$, jadi harapan banyaknya pengambilan kira-kira
$\eu \approx 2.7$: sediakanlah tiga kali putaran topi.

\textbf{17.} Setiap jabat tangan menyumbang $2$ pada cacah derajat
seluruhnya, sehingga jumlah bilangan jabat tangan semua tamu
bernilai genap. Jumlah \omterm{def:g10:numbers:sets}{bilangan bulat} bernilai genap hanya jika banyaknya
suku ganjil itu genap: jadi penjabat ganjil selalu berpasangan. (Pada pesta
bertiga: profil jabat tangan yang mungkin tak pernah punya tepat satu atau
tiga isian ganjil --- periksalah keempat graf yang mungkin.)

\textbf{18.} Untuk $n = 1$: $1 = 1$. Jika
$1^3 + \dots + n^3 = \left(\frac{n(n+1)}{2}\right)^2$, maka dengan menambahkan
$(n+1)^3$:
\[
\frac{n^2(n+1)^2}{4} + (n+1)^3
= \frac{(n+1)^2\left(n^2 + 4n + 4\right)}{4}
= \left(\frac{(n+1)(n+2)}{2}\right)^{\!2} :
\]
itulah pewarisannya. Untuk $n = 3$: $1 + 8 + 27 = 36 = 6^2$.

\textbf{19.} Sebuah lintasan menuju $(n, n)$ menempuh $2n$ langkah dan memotong
diagonal lawannya $x + y = n$ di tepat satu titik kekisi
$(j, n - j)$; separuh pertamanya adalah lintasan dengan $j$ huruf R di antara $n$
langkah ($\binom nj$ pilihan), dan separuh keduanya, yang dibaca mundur,
juga demikian ($\binom nj$ lagi, menurut kesetangkupannya). Dengan menjumlahkan
atas titik potongnya: $\binom{2n}{n} = \sum_j \binom nj^2$ ---
identitas Vandermonde, yang tergambar.

\textbf{20.} Kalikan tahapannya, jumlahkan kasusnya: pelat nomor dan tangan
poker. Sandikan: lintasan sebagai kata RU, pesanan sebagai bintang dan batang.
Cacahlah dua kali: panitia-beserta-ketua, jabat tangan, lintasan yang terpotong
di tengah. Kurangkan lalu perbaiki: topi yang tertukar semua, dengan $\frac1\eu$
sebagai sisanya. Persinggahan berikutnya: cacahan ini di bawah pecahan
peluang, dan pangkat matriks ketetanggaan yang mencacah lintasan
dua bab lagi.

\end{solution}
